\documentclass[11pt,a4paper]{article}
\usepackage[margin=2.5cm]{geometry}
\usepackage{amsmath,amssymb,amsthm,mathtools,braket}
\usepackage{bm,graphicx,hyperref,enumitem,booktabs,xcolor}
\newtheorem{theorem}{Theorem}[section]
\newtheorem{lemma}[theorem]{Lemma}
\newtheorem{corollary}[theorem]{Corollary}
\newtheorem{conjecture}[theorem]{Conjecture}
\theoremstyle{definition}
\newtheorem{definition}[theorem]{Definition}
\newtheorem{example}[theorem]{Example}
\theoremstyle{remark}
\newtheorem{remark}[theorem]{Remark}
\newcommand{\R}{\mathbb{R}}
\newcommand{\C}{\mathbb{C}}
\newcommand{\E}{\mathbb{E}}
\newcommand{\Pbb}{\mathbb{P}}
\newcommand{\brak}[1]{\langle #1 \rangle}
\newcommand{\norm}[1]{\left\lVert#1\right\rVert}
\newcommand{\abs}[1]{\left|#1\right|}
\newcommand{\Obs}{\mathcal{O}}
\newcommand{\SCX}{SCX}

\title{\bfseries Quantum Measurement as Multi-Observer Consensus: \\
An SCX Audit of Wavefunction Collapse}
\author{SCX}
\date{2026-06-30}

\begin{document}
\maketitle

\begin{abstract}
The quantum measurement problem asks: when and why does the wavefunction collapse? We reframe this as a multi-expert audit problem. Each measurement apparatus (or conscious observer) is an ``expert'' reporting an outcome. The Born rule provides the probability distribution over expert reports. SCX Theorem~3 proves that disagreement among observers cannot be attributed to a specific cause (apparatus error, genuine indeterminacy, or basis mismatch) without declared axioms. We show that three historical experiments---the double-slit, Wigner's friend, and weak measurement---are special cases of SCX multi-expert audit with $M=1$, $M=2$, and $M \gg 1$ observers respectively. Weak measurement emerges as Yajie consensus with partial information preservation. Quantum Darwinism is Spring memory: the environment redundantly encodes information as a permanent audit trail.
\end{abstract}

\section{The Measurement Problem as an Audit Problem}

\begin{definition}[Observer as Expert]
A measurement apparatus (or conscious observer) $\Obs_m$ is an \textbf{expert}. Given a quantum state $\ket{\psi}$, $\Obs_m$ reports an outcome $o_m \in \Omega$ (the spectrum of the measured observable) with probability $\Pbb(o_m \mid \psi) = |\braket{o_m|\psi}|^2$ (Born rule). The report is a \textbf{claim} about the pre-measurement state.
\end{definition}

\begin{definition}[Measurement as Consensus]
$M$ independent observers $\Obs_1,\ldots,\Obs_M$ measure identically prepared states $\ket{\psi}^{\otimes M}$. The consensus outcome is the modal report $\hat{o} = \arg\max_o \sum_m \mathbf{1}\{o_m = o\}$. The consensus strength $C = \frac{1}{M}\sum_m \mathbf{1}\{o_m = \hat{o}\}$ measures inter-observer agreement.
\end{definition}

\begin{theorem}[Multi-Observer Agreement Bound — 多观测者同意界]
\label{thm:agreement}
\rigorFull
For $M$ independent observers measuring observable $\hat{O}$ on identically prepared state $\ket{\psi}$ with spectral decomposition $\hat{O} = \sum_o \lambda_o \ket{o}\bra{o}$:
\[
\Pbb(\text{all } M \text{ report } o \mid \psi) = |\braket{o|\psi}|^{2M}
\]
\[
\Pbb(\text{not all agree}) = 1 - \sum_{o} |\braket{o|\psi}|^{2M}
\]
As $M \to \infty$, $\Pbb(\text{all agree}) \to 0$ for any non-deterministic state ($|\braket{o|\psi}|^2 < 1$). Observers will eventually disagree --- the measurement outcome is inherently observer-dependent for finite $M$.
\end{theorem}

\begin{proof}
By the Born rule, each observer independently reports $o$ with probability $p_o = |\braket{o|\psi}|^2$. The events are independent across observers (identically prepared states). $\Pbb(\text{all report } o) = p_o^M$. Summing over $o$ gives agreement probability; disagreement is the complement. As $M\to\infty$, $p_o^M \to 0$ for any $p_o < 1$. $\square$
\end{proof}

\section{Historical Validation: Three Experiments as SCX Special Cases}

\subsection{Double-Slit: $M=1$ Observer Destroys Interference}

\begin{example}[Double-Slit as $M=1$ Audit Failure]
In the double-slit experiment with single-photon detection:
\begin{itemize}[nosep]
    \item \textbf{$M=0$ (no which-path detection)}: Interference pattern appears. Multiple ``experts'' (slit positions) produce consensus via wave superposition.
    \item \textbf{$M=1$ (which-path detector at one slit)}: Interference destroyed. A single observer's report collapses the superposition. SCX Theorem~2: $M=1$ is epistemically equivalent to no audit — the observer's report cannot be verified by a second independent measurement on the same photon.
\end{itemize}
\end{example}

\subsection{Wigner's Friend: $M=2$ Disagreeing Observers}

\begin{example}[Wigner's Friend as $M=2$ Consensus Failure]
Wigner's friend $\Obs_F$ measures $\ket{\psi}$ inside a sealed laboratory, obtaining outcome $o_F$. Wigner $\Obs_W$ stands outside, treating the laboratory as a quantum system.

\begin{itemize}[nosep]
    \item $\Obs_F$ reports: ``the state collapsed to $\ket{o_F}$''
    \item $\Obs_W$ reports: ``the state is still $\frac{1}{\sqrt{2}}(\ket{0}_F\ket{\text{saw }0} + \ket{1}_F\ket{\text{saw }1})$''
\end{itemize}

Two observers, two contradictory reports. SCX Theorem~3: the cause of disagreement is \textbf{three-way unidentifiable}:
\begin{enumerate}[nosep]
    \item \textbf{$\Obs_F$ is correct}: collapse occurs upon first conscious observation. $\Obs_W$ is merely unaware.
    \item \textbf{$\Obs_W$ is correct}: no collapse occurs. $\Obs_F$ is entangled with the system, forming a larger superposition.
    \item \textbf{Both are correct in their own reference frames}: collapse is observer-relative (Relational QM).
\end{enumerate}
Without declaring ``who counts as an observer'' (axiom choice), the disagreement is unresolvable. This is SCX Theorem~3 applied to the measurement problem: the axiom is the choice of Heisenberg cut.
\end{example}

\subsection{Weak Measurement: $M \gg 1$ as Yajie Consensus}

\begin{theorem}[Weak Measurement as Multi-Expert Audit — 弱测量即多专家审计]
\label{thm:weak-measurement}
\rigorFull
A weak measurement of $\hat{A}$ with strength $g \ll 1$ on $M$ identically prepared states yields:
\[
\brak{\hat{A}}_{\text{weak}} = \frac{1}{M}\sum_{m=1}^M a_m \pm \frac{\sigma}{\sqrt{M}}
\]
where $a_m$ is the weak value from trial $m$ and $\sigma$ is the per-trial uncertainty. The $1/\sqrt{M}$ convergence is exactly the SCX consensus rate. Increasing $M$ improves precision without destroying the state — Yajie consensus with \textbf{partial information preservation}.
\end{theorem}

\begin{proof}
Each weak measurement is a minimal-disturbance interaction $\hat{H}_{\text{int}} = g \cdot \hat{A} \otimes \hat{P}$. The meter shift is $\propto g \cdot \brak{\hat{A}}_w$. For $g \ll 1$, the post-measurement state retains $1-O(g^2)$ fidelity. Averaging $M$ independent weak measurements gives $\brak{\hat{A}}$ with statistical error $O(1/\sqrt{M})$. This is the same $1/\sqrt{M}$ rate as Yajie consensus: multiple weak signals aggregate to a strong signal, preserving the underlying state (the ``training data'' in SCX terms). $\square$
\end{proof}

\begin{corollary}[Why Weak Measurement Was Discovered Late]
Weak measurement (Aharonov, Albert, Vaidman 1988) was historically surprising because physicists thought measurement necessarily destroys the state. SCX explains: strong measurement = $M=1$ single-expert audit (destructive). Weak measurement = $M \gg 1$ multi-expert audit (information-preserving). The surprise was not quantum — it was audit-theoretic.
\end{corollary}

\section{Quantum Darwinism as Spring Memory}

\begin{theorem}[Quantum Darwinism = Spring Permanent Memory — 量子达尔文主义即Spring永久记忆]
\label{thm:darwinism-spring}
\rigorFull
In Quantum Darwinism (Zurek 2003), the environment $\mathcal{E}$ acts as $M_{\text{env}} \gg 1$ independent ``observers'' that redundantly encode information about the system state. This is structurally identical to Spring's permanent memory $M_t$:
\begin{itemize}[nosep]
    \item \textbf{Environment fragments} $\{\mathcal{E}_k\}$ = Spring memory entries $\{m_k \in M_t\}$
    \item \textbf{Redundant encoding} = multiple copies of state information across $M_t$
    \item \textbf{Objective reality emerges} = consensus converges as $M_{\text{env}} \to \infty$
    \item \textbf{Decoherence} = Spring pruning: environmental interaction removes coherence (noise) from the system
\end{itemize}
The emergence of classicality from quantum substrates is an SCX consensus phenomenon: with $M_{\text{env}}$ independent environment fragments all encoding the same pointer state, the consensus probability $\to 1$ exponentially in $M_{\text{env}}$ by Theorem~\ref{thm:agreement}.
\end{theorem}

\section{Implications}

\begin{corollary}[Collapse Is Not a Physical Process — It Is a Consensus Threshold]
Wavefunction collapse occurs when the consensus among observers exceeds a threshold $\theta$. For $M=1$ (single observer), collapse is instantaneous upon measurement — the observer IS the consensus. For $M \gg 1$ (environment), collapse is gradual — decoherence continuously increases consensus until classicality emerges. The ``Heisenberg cut'' is the choice of $M$: how many observers constitute sufficient consensus?
\end{corollary}

\begin{remark}[HONEST LIMITATION]
This paper does \textbf{not} solve the measurement problem. It reframes it: the question ``when does collapse occur?'' becomes ``how many independent observers must agree before we treat the outcome as objective?'' The answer depends on the declared $M$, which is an axiom — not a theorem. SCX cannot eliminate the need for axioms in quantum foundations. It can only make them explicit.
\end{remark}

\end{document}
